% \iffalse meta-comment
%
% hit-final-matter.dtx 是前置、正文与后置的切换
%
% \fi
%
% \section{前置、正文与后置的切换}
%
% \changes{v3.2a}{2026/09/13}{学位论文的骨架宏排列顺序照规范写死，\option{struct} 收
%   摘要、正文各章、符号表、缩略语表、图表公式索引各成部件；文献表调总出口
%   \cs{hit@bib@print:} 按后端分发；右开页走 \cs{hit@clear@odd@page:}
%   （\cs{cleardoublepage} 绑着 \cs{if@twoside}，本硕单面时开关是死的），
%   \option{openright} 与 \option{library} 转发到 \option{layout}/\option{openright}；
%   新版面下不发 \cs{ziju}（会盖掉字符网格算出来的 \cs{CJKglue} 与 \cs{ccwd}）；
%   目录里第一章之前不插空行}
%
% ^^A ======================================================================
% ^^A Module final-mainmatter: 正文与前后文切换（frontmatter/mainmatter/backmatter）（hit-thesis）
% ^^A ======================================================================
%    \begin{macrocode}
%<*final-mainmatter>
\bool_if:NT \g__hit_final_bool
  {
%    \end{macrocode}
% \begin{macro}{\cleardoublepage}
% 对于 \textsl{openright} 选项，必须保证章首页右开，且如果前章末页无内容须
% 清空其页眉页脚。
% 如果\textsl{library}为真，则强制设置\textsl{openright}为真。
% \changes{v2.0.10}{2019/06/25}{添加\textsl{openright}和\textsl{library}逻辑}
% \changes{v3.0.0}{2020/05/20}{将glossaries移动到这里}
%    \begin{macrocode}
%    \end{macrocode}
%
% 两个旧的文档类选项转发到新键。它们只设一个布尔，真正读值的是
% \cs{hit@openright:n}，所以在这里把布尔换算成三态。\option{library} 就是全篇
% 不跳，\option{openright} 就是全篇都跳。旧名表那边已经发过一次更名提示。
%    \begin{macrocode}
\bool_if:NT \g__hit_etitleauto_bool
  { \tl_gset:Nn \g_hit_titlepage_title_en_xiaoer_tl { auto } }
\bool_if:NT \g__hit_openright_bool
  { \tl_gset:Nn \g__hit_openright_tl { true } }
\bool_if:NT \g__hit_library_bool
  { \tl_gset:Nn \g__hit_openright_tl { false } }
\cs_set_eq:NN \hit@clear@double@page \cleardoublepage
\cs_new:Npn \hit@clear@empty@double@page
  { \clearpage { \pagestyle { hit@empty } \hit@clear@double@page } }
\cs_set_eq:NN \cleardoublepage \hit@clear@empty@double@page
%    \end{macrocode}
%
% \cs{hit@clear@odd@page:} 翻到下一个奇数页，\emph{不看} \cs{if@twoside}。
%
% \cs{cleardoublepage} 内核那一版写的是 \cs{if@twoside} 为真才补空白页，而
% \option{oneside}/\option{twoside} 在这个类里是按学位派的（本硕单面、博士双面）。
% 右开页要是走它，本硕写 \option{open-right}\texttt{=true} 也补不出空白页——
% 开关成了摆设，学位耦合又回到了 \option{oneside} 里。
%
% \option{oneside}/\option{twoside} 仍旧照派，它还管着页眉奇偶与页边距，那两样
% 本来就该跟着单双面走。这里只是把\emph{跳页}这一件事从它手里拿出来。
%
% 空白页的页眉照 \cs{hit@clear@empty@double@page} 的办法压掉。
%    \begin{macrocode}
\cs_new_protected:Npn \hit@clear@odd@page:
  {
    \clearpage
    {
      \pagestyle { hit@empty }
      \int_if_odd:nF { \value { page } } { \hbox:n { } \newpage }
    }
  }
%    \end{macrocode}
%
% \end{macro}
% \begin{macro}{\frontmatter}
% 我们的单面和双面模式与常规的不太一样。
%    \begin{macrocode}
\cs_set:Npn \frontmatter
  {
    \tl_gset:Nn \g__hit_matter_now_tl { frontmatter }
    \hit@openright:n { frontmatter }
    \@mainmatterfalse
    \pagenumbering { Roman }
    \pagestyle { hit@empty }
%    \end{macrocode}
%
% 窝工深圳模板范例，一行34个字
% \changes{v3.0.11}{2020/10/29}{窝工深圳模板范例，一行34个字}
% 添加33行字选项
% \changes{v3.0.13}{2020/11/10}{添加33行字选项}
%
% \cs{ziju} 是旧版面把一行凑成 34 个字的手段，靠的是给每个汉字后面加一段
% 与字号成比例的空隙。新版面这活儿由字符网格干，两者不能叠。
%
% 而且不是「叠了偏大」这么简单：\cs{ziju} 底层走 \pkg{ctex} 的
% \cs{ctex\_update\_ccglue:}，那一句是 \cs{cs\_set\_protected:Npn} \cs{CJKglue}，
% 把 \cs{docxlike\_format\_apply\_body:} 设好的网格字距整个换掉，还顺手给 \cs{ccwd}
% 加了 5.58\%。\cs{mainmatter} 不复原，所以一发就是全文。本科那一档实测
% \cs{ccwd} 从 12.50099\,pt 变成 12.71\,pt，字距从规范的 12.4543\,bp
% 变成 12.658\,bp，一行 34 个字累计差 6.9\,bp。
%
% 判据取 \cs{g\_docxlike\_page\_char\_dim}：它非零就说明 \cs{docxlike\_page\_setup:n}
% 走过，字符网格立起来了，与 \file{src/docxlike/docxlike-format.dtx} 里那几处同一个判据。
%    \begin{macrocode}
    \hit@openright:n { frontmatter }
    \dim_compare:nNnT { \g_docxlike_page_char_dim } = { 0pt }
      {
        \bool_if:NTF \g__hit_ziju_word_bool
          { \ziju { 0.05416667 } }
          { \bool_if:NT \g__hit_ziju_regu_bool { \ziju { 0.05580808 } } }
      }
  }
%    \end{macrocode}
%
% \end{macro}
% \begin{macro}{\mainmatter}
% 根据打印店（伪官方）的猛虎式操作，\cs{mainmatter}命令的逻辑是，双面打印时第一章必须在奇数页
% （不看文档别怪我）。
% \changes{v2.0.11}{2018/06/27}{设置第一章必须在奇数页}
%
% 原先目录里第一章之前会空一行，由 \option{tocblank} 控制。规范中没有这一条，
% 整段去掉，\option{tocblank} 写了也不起作用。
%    \begin{macrocode}
\cs_set:Npn \mainmatter
  {
%    \end{macrocode}
%
% 添加英文版限制
% \changes{v3.0.0}{2020/05/21}{添加英文版限制}
%    \begin{macrocode}
    \tl_gset:Nn \g__hit_matter_now_tl { mainmatter }
    \hit@openright:n { mainmatter }
    \@mainmattertrue
    \pagenumbering { arabic }
    \pagestyle { hit@headings }
  }
%    \end{macrocode}
%
% \end{macro}
% \begin{macro}{\backmatter}
%    \begin{macrocode}
\cs_set:Npn \backmatter
  {
    \tl_gset:Nn \g__hit_matter_now_tl { backmatter }
    \hit@openright:n { backmatter }
    \@mainmattertrue
  }
%    \end{macrocode}
% \end{macro}
%
% \begin{macro}{\hitfrontmatter,\hitbackmatter}
% 规范定死的顺序，各部件由 \option{struct} 族声明，没声明的不排。
% 键的声明与存取在 \file{src/engine/hit-structure.dtx}。
%
% \option{cover} 与 \option{abstract} 两个文件排在最前面读，跟别的部件「读到哪里
% 排到哪里」不是一回事。这两个文件本身不排任何东西：封面那个只设元信息，摘要那个
% 里的 \env{abstract-zh} 与 \env{abstract-en} 只把正文收起来。真正排版的是
% \cs{makecover}，它接着封面一起排摘要，所以两个文件都得赶在它之前读到。
%    \begin{macrocode}
\cs_set_protected:Npn \hitfrontmatter
  {
    \bool_gset_true:N \g__hit_matter_front_bool
    \hit@structure@input:n { cover }
    \hit@structure@input:n { abstract }
    \frontmatter
    \makecover
    \hit@structure@input:n { denotation }
%    \end{macrocode}
%
% 符号表取文件名，内容是手写的排版材料；缩略语表是三值，那一页由类从
% \option{abbr} 族的词条生成，没有文件可读。第三个部件把两者排在同一页，
% 标题另取一对常量，两段各带默认小标题。
%    \begin{macrocode}
    \hit@structure@input:n { list-of-symbols }
    \hit@structure@switch:nnn { list-of-abbreviations }
      { \bool_if_p:N \g__hit_en_bool } { \listofabbreviations }
    \hit@structure@merged@symbols:
    % 目录:规范要求都有,auto 即排
    \hit@structure@switch:nnn { table-of-contents } { \c_true_bool }
      { \tableofcontents }
%    \end{macrocode}
%
% 三个索引。从前这三样(连同缩略语表)写死在 \cs{tableofcontents} 内部，条件是
% \option{lang}\texttt{=en}，于是「要目录」等于连带要索引，中文版想要图索引也
% 要不到。现在各是一个部件，取值 \texttt{auto}(不写即是)时按规范默认走——规范
% 只对英文版有要求，所以默认条件是「这本是不是英文版」。
%    \begin{macrocode}
    \hit@structure@switch:nnn { list-of-tables }
      { \bool_if_p:N \g__hit_en_bool } { \listoftables }
    \hit@structure@switch:nnn { list-of-figures }
      { \bool_if_p:N \g__hit_en_bool } { \listoffigures }
    \hit@structure@switch:nnn { list-of-equations }
      { \c_false_bool } { \listofequations }
    \hit@structure@here:nnn { acknowledgements } { front }
      { \exp_args:Nv \include { g__hit_structure_acknowledgements_tl } }
  }
%    \end{macrocode}
%
% \begin{macro}{\hitmainmatter}
% \cs{hitmainmatter} 起正文，再按 \option{mainmatter} 逐个 \cs{include}。
% \option{mainmatter} 是唯一取一串文件名的部件，其余都是一个部件一个文件。
%
% 不写 \option{mainmatter} 也行，\cs{hitmainmatter} 照样起正文，章节自己 \cs{include}，
% 跟以前一样。
%    \begin{macrocode}
\cs_set_protected:Npn \hitmainmatter
  {
    \bool_gset_true:N \g__hit_matter_main_bool
    \mainmatter
    \hit@structure@include@list:n { mainmatter }
  }
%    \end{macrocode}
% \end{macro}
%
%    \begin{macrocode}
\cs_set_protected:Npn \hitbackmatter
  {
    \bool_gset_true:N \g__hit_matter_back_bool
    \backmatter
    \hit@structure@part:nn { bibliography } { \hit@bib@print: }
    \hit@structure@part:nn { appendix }
      {
        \begin { appendix }
          \hit@structure@input:n { appendix }
        \end { appendix }
      }
    \hit@structure@include:n { publications }
    % 答辩决议:规范只对博士要求,auto 即「是博士才排」。原先这条三态是手写在
    % 这里的,现在收进 \cs{hit@structure@scanned:nnnn} 统一处理
    \hit@structure@scanned:nnnn { defense } { hit@resolution@place }
      { \bool_if_p:N \g__hit_doctor_bool } { \hit@resolution@place }
    \hit@structure@include:n { index }
    % 原创性声明与使用权限:规范要求都有,auto 即排
    \hit@structure@scanned:nnnn { declarations } { hit@declarations@place }
      { \c_true_bool } { \hit@declarations@place }
    \hit@structure@here:nnn { acknowledgements } { back }
      { \exp_args:Nv \include { g__hit_structure_acknowledgements_tl } }
    \hit@structure@include:n { resume }
  }
%    \end{macrocode}
% \end{macro}
%    \begin{macrocode}
  }
%</final-mainmatter>
%    \end{macrocode}
%
% \Finale
